% AIPL 3実装の同期 ---- 実装レポート
%   lualatex -interaction=nonstopmode aipl_three_impl_parity_ja.tex   （2回）
\documentclass[11pt,a4paper]{ltjsarticle}
\usepackage{amsmath,amssymb}
\usepackage{listings}
\usepackage{xcolor}
\usepackage{geometry}
\usepackage{array}
\usepackage{booktabs}
\usepackage{longtable}
\usepackage[hidelinks]{hyperref}
\geometry{margin=22mm}

\definecolor{codebg}{gray}{0.96}
\definecolor{okgreen}{rgb}{0.0,0.45,0.15}
\definecolor{badred}{rgb}{0.65,0.10,0.10}
\lstset{
  basicstyle=\ttfamily\small, backgroundcolor=\color{codebg},
  frame=single, framesep=4pt, rulecolor=\color{gray},
  breaklines=true, showstringspaces=false,
  columns=fullflexible, keepspaces=true,
  keywordstyle=\bfseries, commentstyle=\itshape\color{gray!130},
}
\lstdefinelanguage{AIPL}{
  morekeywords={class,method,var,new,now,future,await,reply,send,become,
    timeout,else,if,while,do,select,case,self,sender,remote,print,return,
    int,float,string,bool,unit,any,true,false},
  sensitive=true, morecomment=[l]{//}, morestring=[b]",
}
\newcommand{\ok}{\textcolor{okgreen}{\textbf{○}}}
\newcommand{\bad}{\textcolor{badred}{\textbf{×}}}

\title{AIPL 3実装の同期\\
{\large ---\ 拡張子の統一、仕様追随、そして調べる過程で出た6つの不具合 ---}}
\author{児玉工房 / AICE $\cdot$ AIPL プロジェクト}
\date{2026年08月01日 00:25}

\begin{document}
\maketitle

\begin{abstract}
AIPL には実装が3つある ---- 正である OCaml 版、Python の CLI 実装（Py-I）、
ブラウザで動く JavaScript 実装（JS-I）。OCaml 版に加えた一連の仕様変更に
残る2つが追いついておらず、\textbf{ガイドのサンプル8本が Py-I でも JS-I でも
1本も動かない}状態になっていた。

本レポートはこれを揃えた作業の記録である。やったことは二つ。
ソースの拡張子を \texttt{.abcl} から \texttt{.aipl} へ統一し（4リポジトリ 478 本）、
Py-I と JS-I を最新の言語仕様へ追随させた。結果、サンプルは3実装すべてで動き、
出力も一致する。

ただし本レポートの主眼は作業一覧ではない。\textbf{揃える過程で6つの不具合が
出てきたこと}と、\textbf{それがどうやって見つかったか}である。うち3つは
正であるはずの OCaml 版のバグだった。とくに \texttt{g4\_select} の挙動差を
「スケジューリングの違い」と片付けずに\emph{どちらが正しいか}を調べたことが、
トップレベル送信の順序破壊 $\to$ select の返信経路の欠落 $\to$ アクターを
引数で渡せない、という連鎖の発見につながった。

見つけ方の型は3つに集約できる。\textbf{3実装を並べて件数を比べる}、
\textbf{機械的な置換をせず構文木で判定する}、\textbf{検査の指摘は件数でなく
実物を開いて確かめる}。最後の型は実際に効き、reply 検査の 62 件の指摘のうち
59 件が誤検出であることを暴いた。
\end{abstract}

\tableofcontents

\section{三つの実装}

\begin{center}
\small
\begin{tabular}{@{}llp{6.2cm}@{}}
\toprule
 & 場所 & 備考 \\
\midrule
OCaml 版 & \texttt{\textasciitilde/aios/abclcp} & \textbf{正}。10.8k 行。
  REPL から \texttt{load} / \texttt{compile} \\
Py-I & \texttt{\textasciitilde/test-bed/aios-claude/src/python-aipl} &
  CLI。15.7k 行。Lark 文法（\texttt{grammar.lark}）\\
JS-I & \texttt{\textasciitilde/projects/drone-hil/abcl} &
  \textbf{ブラウザ実装}。5.2k 行。jison 文法を
  \texttt{parser.js} へ生成 \\
\bottomrule
\end{tabular}
\end{center}

JS-I がブラウザ実装であることは以降の随所に効いてくる。デモは HTML ページで、
\texttt{index.html} が \texttt{parser/parser.js} を \texttt{<script>} で読み込み、
\texttt{interpreter.js} を ES モジュールとして import する。検証のために
node から \texttt{parser.yy} に \texttt{ast.js} を束ねて \texttt{Interpreter} を
回すヘッドレスの足場を組んだ。

\subsection{着手前の状態}

ガイドのサンプル（\texttt{docs/samples/guide/g1..g8}）を各実装に流したところ、
\textbf{Py-I も JS-I も8本すべて失敗}した。Py-I の実際のエラーはこうである。

\begin{lstlisting}
g8: [parse error] No terminal matches 'f' ... var flag = false;
g7: [parse error] Unexpected token Token('COLON', ':')
g1: [parse error] Unexpected token Token('PLUS', '+')
\end{lstlisting}

\noindent 真偽値リテラル、戻り値型注釈 \verb|: T|、文字列連結 \verb|++| の
いずれもが無い。JS-I も \texttt{grammar.jison} に \verb|++|・\verb|:|・
\verb|!{| の規則そのものが無かった。

\section{拡張子の統一}

言語名は AIPL なのにソースの拡張子が \texttt{.abcl} のままだった。前身の
ABCL に由来する名残で、3実装をまたいで混乱の元になっていた。

\begin{center}
\small
\begin{tabular}{@{}lrrl@{}}
\toprule
リポジトリ & 改名 & 参照追随 & コミット \\
\midrule
\texttt{aios/abclcp} & 93 & 50 & \texttt{cccb29d} \\
\texttt{test-bed/aios-claude} & 380 & 443 & \texttt{2ac81b9} \\
\texttt{projects/drone-hil} & 3 & 6 & \texttt{b78ec34} \\
\texttt{aipl\_line\_simulator} & 2 & 10 & \texttt{7856575} \\
\midrule
合計 & \textbf{478} & \textbf{509} & \\
\bottomrule
\end{tabular}
\end{center}

履歴を保つため \texttt{mv} ではなく \texttt{git mv} を使った。参照は
\texttt{.sh}・\texttt{.bat}・\texttt{.md}・\texttt{.py}・\texttt{.json}・
\texttt{.tex}・\texttt{.js} など広範に及んだが、拡張子を決め打ちしている
\emph{コード}は \texttt{aipl\_dashboard.py} の2箇所（\texttt{endswith(".abcl")}）
だけだった。処理系のパーサ・字句解析器・REPL は拡張子を見ていない。

検証は「型検査の結果が改名前と1件ずつ一致するか」で行い、
\textbf{93本中93本が完全一致}した。うち27本は \texttt{PARSE\_ERROR} だが、
これは改名前から同じ旧構文の残骸であり本件の退行ではない。

\section{仕様の追随}

OCaml 版に入っていて Py-I / JS-I に無かったものを移した。

\begin{center}
\small
\begin{tabular}{@{}lccc@{}}
\toprule
 & OCaml & Py-I & JS-I \\
\midrule
\verb|++|（文字列連結）& 既 & \ok 追加 & \ok 追加 \\
真偽値リテラル & 既 & \ok 追加 & \ok 追加 \\
優先順位の分離 & 既 & \ok 追加 & \ok 追加 \\
\texttt{timeout ... else}（期限）& 既 & \ok 追加 & \ok 追加 \\
戻り値型注釈 \verb|: T| & 既 & \ok 別名追加 & \ok 追加 \\
効果注釈 \verb|!{...}| & 既 & 既 & \ok 追加 \\
\verb|+| を数値専用に & 既 & \ok 移行 & \ok 追加 \\
引数の型注釈 \verb|method m(x: T)| & \ok 追加 & 既 & \ok 追加 \\
reply 線形性の検査 & 既 & \ok 追加 & \ok 追加 \\
期限の警告 & 既 & \ok 追加 & \ok 追加 \\
\bottomrule
\end{tabular}
\end{center}

優先順位は OCaml 版に合わせて分離した。低い順に

\[
\text{等値} < \text{比較} < \text{\texttt{++}} < \text{加減} < \text{乗除}
< \text{単項マイナス}
\]

\noindent であり、\textbf{二項演算子はすべて左結合}である（等値も含む）。
これは実測で確かめた ---- \verb|1 == 2 == false| は
\verb|((1 == 2) == false)| と読まれて \texttt{true} になり、
\verb|- - 3| は \texttt{3} になる。

\subsection{Py-I は書式が違っていただけの箇所があった}

Py-I の \texttt{grammar.lark} には既に \texttt{return\_anno} と
\texttt{effect\_anno} があった。無かったのではなく\textbf{書式が違って}いた。

\begin{center}
\small
\begin{tabular}{@{}lll@{}}
\toprule
 & OCaml 版 & Py-I \\
\midrule
戻り値型注釈 & \verb|method m(x) : T| & \verb|method m(x) -> T| \\
\bottomrule
\end{tabular}
\end{center}

両方受けるようにした。注釈が黙って捨てられていないことは、パースした
AST を直接覗いて確かめた（\texttt{ret='int' eff=['log']}）。
「通ったが無視されていた」は AIPL の道具立てで何度も踏んだ罠である。

\section{\texttt{+} から \texttt{++} への移行 ---- 正規表現を使わなかった理由}
\label{sec:migrate}

OCaml 版では \verb|+| は数値専用で、文字列連結は \verb|++| に分離済みである
（オーバーロードの曖昧さを潰すため）。同じことを Py-I でやったところ、
テストが 15通過/5失敗 $\to$ 6通過/14失敗 に悪化した。調べると
\textbf{既存の \texttt{.aipl} 資産 437 本のうち 263 本が \texttt{+} で文字列を
連結}していた。

移行が必要だが、\textbf{正規表現で置換してはいけない}。\verb|+| は左結合で
\verb|++| より強く結合するので、連鎖の \verb|+| を1つだけ替えると意味が変わる。

\begin{lstlisting}[language=AIPL]
"a" + b + c        // (("a" + b) + c)  ---- 全体が文字列連結
"a" ++ b + c       // "a" ++ (b + c)   ---- b と c が数値なら加算に化ける
\end{lstlisting}

そこで構文木で判定した。Lark を \texttt{propagate\_positions=True} で
生の木に落とし、\texttt{add\_expr} の\emph{連鎖ごと}に
「オペランドのどれかが文字列か」を見て、該当する連鎖の \verb|+| トークンの
位置をすべて集め、\textbf{末尾から}置換する（\verb|+| 1文字 $\to$
\verb|++| 2文字で以降の位置がずれるため）。

\begin{center}
\small
238 本 / 5448 箇所を書き換え
\end{center}

\subsection{構文木は文字列リテラルの中を見ない}

これで済んだと思ったらテストが2本落ちた（\texttt{\_test\_dynamic},
\texttt{\_test\_methodpatch}）。原因は、実行時にコンパイルする AIPL ソースを
\emph{文字列として}持つファイルがあったこと。

\begin{lstlisting}[language=AIPL]
"    print(\"  [Counter] tick -> n=\" + n);\n" ++
\end{lstlisting}

\noindent 外側の連結は \verb|++| に直っているのに、文字列の中の \verb|+| は
構文木からは見えないので残る。\verb|\"| を含む行を別途処理して解消した
（8本11行）。

\subsection{検証}

移行後、437本中413本がパースできる。失敗24本は移行前と同数・同一である
（\texttt{Generics.aipl} の型引数など、元から通らない旧構文）。テストは
15通過/5失敗で\textbf{基準と完全一致}し、新規失敗はゼロだった。

\section{見つかった6つの不具合}

ここからが本題である。揃える作業そのものより、\textbf{揃えようとして初めて
見えた不具合}のほうが多く、しかもその半分は正であるはずの OCaml 版にあった。

\subsection{(1) トップレベルの \texttt{send} がプログラム順を破っていた}
\label{sec:order}

\texttt{g4\_select} の出力が3実装で食い違っていた。

\begin{center}
\small
\begin{tabular}{@{}ll@{}}
\toprule
OCaml 版 & \texttt{direct:1 / waiting / timed out} \\
Py-I, JS-I & \texttt{waiting / got 1} \\
\bottomrule
\end{tabular}
\end{center}

これを「スケジューリングの違い」で片付けず、\textbf{どちらが正しいのかを
最小の例で測った}。

\begin{lstlisting}[language=AIPL]
class A {
  method first()  : unit !{log} { print("1st: first"); }
  method second() : int  !{log} { print("2nd: second"); reply(7); }
}
var a = new A();
send a.first();
print("now returned " ++ (now a.second() timeout 1000 else 0));
\end{lstlisting}

\begin{center}
\small
\begin{tabular}{@{}ll@{}}
\toprule
OCaml 版 & \texttt{2nd: second / now returned 7 / 1st: first} \\
Py-I, JS-I & \texttt{1st: first / 2nd: second / now returned 7} \\
\bottomrule
\end{tabular}
\end{center}

\textbf{OCaml 版だけが FIFO を破っていた}。5回ずつ走らせて、どちらも
決定的であることを確かめている（競合ではない）。

原因はメールボックスではなかった。アクターの \texttt{queue} は単一の
\texttt{Queue.t} で、\texttt{actor\_loop} も \texttt{Queue.pop} の厳密な FIFO、
\texttt{send\_message} も mutex 下で即 \texttt{push} する。つまり配送は正しい。
ずれていたのは \texttt{repl\_thread.ml} の \texttt{compile} で、
\textbf{トップレベルの \texttt{Send} だけを \texttt{pending\_global\_sends} に
積んで全文の実行後にまとめて送っていた}のに、\texttt{CallStmt}
（\texttt{print(now ...)} など）はその場で実行していたことだった。

後回しは前方参照のために入っていたと思われるが、対象93本を調べたところ
\texttt{send} が \texttt{var ... = new ...} より前に来るファイルは\textbf{0本}
だった。それでも壊さないよう、\textbf{宛先が既に生成されていればその場で送り、
まだ無い場合だけ後回しにする}条件付きに変えた（\texttt{df66dca}）。

これで \texttt{g4\_select} は
\texttt{waiting / got 1 / via select:1} となり、ガイドの説明
「select で job を待つ。case 本体の reply は job への返信になる」と一致した。

\subsection{(2)(3) select の case 本体の reply が呼び出し元へ戻らない}

(1) を直したことで \texttt{g4} が\emph{実際に} select 経路を通るようになり、
そこで Py-I と JS-I が破綻していることが分かった。

\begin{center}
\small
\begin{tabular}{@{}ll@{}}
\toprule
Py-I & 値が返らず \texttt{[REPLY] via select:1} とログするだけ \\
JS-I & \texttt{now deadlock: slot rs-1 not fulfilled} で\textbf{失敗} \\
\bottomrule
\end{tabular}
\end{center}

原因は両方とも同じで、case 本体を実行するときに\textbf{受け取ったメッセージの
返信先を据えていなかった}こと。親フレーム（select を書いたメソッド自身）の
返信先を辿ってしまっていた。Py-I は \texttt{drained[i]} の
\texttt{reply\_future} を捨てており、JS-I は
\texttt{actor.\_\_currentSlotId} を差し替えていなかった。それぞれ据えるよう
直した（\texttt{e7b6295}, \texttt{c2929f5}）。

\subsection{(4) 期限検査がトップレベルを見ていなかった ---- 両実装で}

JS-I に期限検査を移し、OCaml 版と\textbf{件数を並べて}比べたところ合わなかった。

\begin{center}
\small
\begin{tabular}{@{}lccc@{}}
\toprule
期限警告の件数 & OCaml & Py-I & JS-I \\
\midrule
\texttt{g3\_actors} & 3 & 1 & 1 \\
\texttt{g6\_effects} & 3 & 1 & 1 \\
\bottomrule
\end{tabular}
\end{center}

原因はクラスのメソッドしか走査しておらず、
\texttt{print(now front.place(3));} のような\textbf{トップレベルの
\texttt{now}} を見ていなかったこと。JS-I を直して比べ直したら、
\textbf{同じ穴が Py-I にもあった}（\texttt{GlobalStmt} を通っていなかった）。
両方直して全項目が一致した（\texttt{22c0569}, \texttt{9c7696c}）。

\textbf{1実装だけ見ていたら両方見逃していた。}

\subsection{(5) JS-I の \texttt{wait} がブロックしない}

\texttt{g7\_deadline} だけが最後まで揃わなかった。

\begin{center}
\small
\begin{tabular}{@{}ll@{}}
\toprule
OCaml 版, Py-I & \texttt{slow(timed out) = 0} \\
JS-I & \texttt{slow(timed out) = 42}（期限が発火しない）\\
\bottomrule
\end{tabular}
\end{center}

OCaml 版と Py-I の \texttt{wait(ms)} はメソッド実行中にブロックするので、
その後の \texttt{reply} も遅れる。JS-I の \texttt{wait} は
\texttt{actor.\_\_nextDelay} を立てて\emph{次の}メッセージ配送を遅らせるだけで、
実行中の処理はブロックしない。ブラウザで同期ブロックすると画面が固まるための
設計であって手抜きではない。

意味を変えるのは実行基盤に踏み込むので、\textbf{先に影響を数えた} ----
drone-hil 全体で \texttt{wait} と \texttt{reply} を同居させているメソッドは
\textbf{0個}だった（\texttt{wait} はデモのアニメーションの間合いにしか
使われていない）。そこでブロックの代わりに
「この先の \texttt{reply} を ms だけ遅らせる」形で近似した（\texttt{4b6f2b5}）。

\begin{itemize}
\item \texttt{wait}: \texttt{actor.\_\_replyDelayMs} に加算
\item \texttt{Reply}: 遅延があれば即 fulfill せず
      \verb|{slotId, value, notBefore}| を積む
\item \texttt{\_drainOneStep}: 毎回フラッシュし、待ちが残っていれば
      「まだ進む余地がある」として \texttt{true} を返す
      （\texttt{false} だと \texttt{drainUntilSlot} が deadlock 例外を投げる）
\end{itemize}

これで5回とも \texttt{slow(timed out) = 0} で安定した。

\subsection{(6) アクターをメッセージの引数として渡せない}
\label{sec:actorarg}

引数の型注釈を入れたあとに見つかった。

\begin{lstlisting}[language=AIPL]
var d = new D();
send p.take(d);      // take が呼ばれない。エラーも出ない
\end{lstlisting}

アクター同士が相手を教え合う基本的な書き方が封じられていた。注釈の有無に
関わらず再現するので、注釈とは無関係の既存の不具合である。原因は\textbf{3つ
重なっていた}。

\paragraph{(6a) 値を式へ戻すときアクターを文字列に潰していた。}
メッセージ引数は値へ評価してから \texttt{expr\_of\_value} で式に戻して送る。
その \texttt{expr\_of\_value} が

\begin{lstlisting}[language={[Objective]Caml}]
| VActor (n,_) -> String ("<actor:" ^ n ^ ">")
\end{lstlisting}

\noindent としており、受け手に届くのはアクターではなく文字列だった。
真偽値も \texttt{String "true"} に潰していた（\texttt{Bool} の式型ができたのに
古いまま）。\texttt{expr\_desc} に \texttt{ActorRef} を追加し、
\texttt{VActor $\to$ ActorRef}、\texttt{VBool $\to$ Bool} とした。

\paragraph{(6b) トップレベルの \texttt{send} が引数を評価していなかった。}
生の式のまま \texttt{send\_message} に渡していたので、受け手の環境で
\texttt{Var "d"} が解決できず黙って落ちていた
（\texttt{actor\_loop} が例外を握って web イベントに積むだけなので無反応に見える）。

\paragraph{(6c) \texttt{VActor} が実体名を持っていなかった。}
\texttt{VActor} の第1要素はクラス名なのだが、
\texttt{resolve\_actor\_from\_term} はそれを\emph{実体名}として読んでおり、
コード内で一貫していない。実体名とクラス名が一致する場合と、変数名で引ける
場合にだけ動いていた ---- 引数で受け取ったアクターはそのどちらでもない。

\texttt{actor\_value \textasciitilde name \textasciitilde cls} を用意し、
第1要素はクラス名のまま実体名を表の \texttt{"\_\_actor\_name"} に入れ、
\texttt{actual\_local\_target} と \texttt{resolve\_actor\_from\_term} が
「実体名 $\to$ 第1要素 $\to$ 変数名」の順で引くようにした。
第1要素の意味は変えていないので \texttt{typeof} 等は無傷である
（\texttt{cbaa7bc}）。

\section{見つけ方の型}

同じ発見が3つの型に集約できる。他の多実装プロジェクトにもそのまま効くと思う。

\subsection{3つ並べて件数を比べる}

\S 5.4 がそれである。1実装を見て「出た」で満足すると過小報告に気づかない。
同じ入力に対する\emph{件数}を並べると、単体では正しく見える実装の穴が出る。

\subsection{機械的な置換をせず構文木で判定する}

\S\ref{sec:migrate} がそれである。演算子の優先順位が絡む書き換えは、
正規表現では\emph{意味を変えずに}行えない。構文木で連鎖を判定し、
バイト位置で末尾から置換する。ただし\textbf{構文木は文字列リテラルの中を
見ない}ので、実行時コンパイルを持つ言語ではそこが漏れる。

\subsection{検査の指摘は件数でなく実物を開いて確かめる}

reply の線形性検査を入れたとき、全437本で\textbf{28本62件}の指摘が出た。
件数だけ見れば「よく効いている」と読めるが、開いてみると全て

\begin{lstlisting}[language=AIPL]
method add(x: int, y: int) -> int { return x + y; }
\end{lstlisting}

\noindent の形だった。\textbf{Py-I のメソッドは OCaml 版と違い
\texttt{return} でも返せる}。\texttt{return} を勘定に入れたら
\textbf{62件が3件に落ちた}。その前段でも、\texttt{function}
（\texttt{reply} ではなく \texttt{return} で返す）に検査をかけて
\texttt{\_test\_linear} を落としている。

残った3件は\textbf{本物}だった。

\begin{lstlisting}[language=AIPL]
method read() {
  if (kind == "Formula") { var v = eval_one(...); reply(v); }
  reply(val);            // if に else も脱出も無いので落ちてくる
}
\end{lstlisting}

\noindent \texttt{kind == "Formula"} のとき2回 \texttt{reply} する。
\texttt{Future.set} が冪等なので実害は出ていないが（先勝ちで \texttt{v} が
返る）、2つ目は死んだコードである。

\begin{quote}
\small
指摘は「出たら本物」でないと使われなくなる。件数で納得すると、
誤検出だらけの検査を仕込むことになる。
\end{quote}

\subsection{おまけ ---- 警告が出力経路に載っているか}

Py-I に期限の警告を足したのに、\texttt{--type-check} にも
\texttt{type\_check()} にも現れなかった。module 直下の \texttt{check()} が
issues しか返していなかったためである。\textbf{検査を足しただけで見えない
ままなら、無いのと同じ}である。

\section{退行させた一度と、その気づき方}

引数の型注釈を JS-I に入れたとき、\texttt{g4\_select} が
\texttt{now deadlock: slot rs-1 not fulfilled} に戻った。
\texttt{params} を使う文法規則が \texttt{method\_decl} と
\texttt{select\_case} の\textbf{2つ}あり、片方だけ直したためである。
\texttt{params} を \verb|{name, ty}| の並びに変えた結果、case 本体の束縛
\texttt{localEnv[p]} がオブジェクトをキーにしてしまい、引数が見えなくなっていた。

\textbf{ガイドの8本を毎回流していなければ気づかなかった。}
検証を「変更に関係ありそうな範囲」に絞らないことが効いた場面である。

\section{最終状態}

\subsection{ガイドのサンプル}

\begin{center}
\small
\begin{tabular}{@{}lccc l@{}}
\toprule
 & OCaml & Py-I & JS-I & 一致 \\
\midrule
\texttt{g1\_hello} & \ok & \ok & \ok & 一致 \\
\texttt{g2\_now\_future} & \ok & \ok & \ok & 一致 \\
\texttt{g3\_actors} & \ok & \ok & \ok & 一致 \\
\texttt{g4\_select} & \ok & \ok & \ok & 一致 \\
\texttt{g6\_effects} & \ok & \ok & \ok & 一致 \\
\texttt{g7\_deadline} & \ok & \ok & \ok & 一致 \\
\texttt{g8\_bool\_equality} & \ok & \ok & \ok & 一致 \\
\texttt{g9\_actor\_arg} & \ok & \ok & \ok & 一致 \\
\bottomrule
\end{tabular}
\end{center}

\noindent（\texttt{g5\_web} は Web サーバを立てるので出力比較の対象外）
着手前は Py-I・JS-I とも \textbf{0/8} だった。

\subsection{\texttt{g9}: アクターを引数で渡す}

本作業で追加したサンプルである。\S\ref{sec:actorarg} の修正がそのまま効く。

\lstinputlisting[language=AIPL]{samples/guide/g9_actor_arg.aipl}

\begin{center}
\small 3実装とも \texttt{got actor / after send / pong}
\end{center}

なお最初は \texttt{typeof(x)} を使って書いていたが、\textbf{\texttt{typeof} は
JS-I に無い}（\texttt{Unknown function: typeof}）ため3実装共通で動かず、
外した。サンプルは3実装で動くことに意味がある。

\subsection{残っている差}

\begin{center}
\small
\begin{tabular}{@{}lp{9.0cm}@{}}
\toprule
項目 & 内容 \\
\midrule
\texttt{typeof} の書式 &
  OCaml: \verb|actor(D) { ping : () -> unit }|、
  Py-I: \verb|actor(D, methods=[ping])|。JS-I には存在しない \\
メソッド内の \texttt{return} &
  Py-I は受けるが JS-I は受けない。
  \texttt{bench/*.aipl} 3本がこれで止まっている \\
アクター型の区別 &
  OCaml の \texttt{unify} は \texttt{TActor} 同士を無条件に成功させる。
  \texttt{actor[A]} と \texttt{actor[B]} が単一化できてしまう \\
出力のフラッシュ &
  OCaml の REPL がスクリプト終了時に最後の1行を取りこぼすことがある。
  実装差ではなく出力の揺らぎ \\
\bottomrule
\end{tabular}
\end{center}

\section{コミット一覧}

\begin{center}
\small
\begin{longtable}{@{}lll@{}}
\toprule
リポジトリ & コミット & 内容 \\
\midrule
\endhead
abclcp & \texttt{cccb29d} & 拡張子を \texttt{.aipl} へ（93本）\\
abclcp & \texttt{df66dca} & トップレベル \texttt{send} の順序（\S\ref{sec:order}）\\
abclcp & \texttt{068bdcd} & 引数の型注釈 \verb|method m(x: T)| \\
abclcp & \texttt{cbaa7bc} & アクターを引数で渡せるように（\S\ref{sec:actorarg}）\\
abclcp & \texttt{47e3fa1} & \texttt{g9} から \texttt{typeof} を外す \\
\midrule
aios-claude & \texttt{2ac81b9} & 拡張子を \texttt{.aipl} へ（380本）\\
aios-claude & \texttt{ce60fbd} & Py-I を最新仕様へ追随 \\
aios-claude & \texttt{9ccc052} & \verb|+| を数値専用に、241本を移行 \\
aios-claude & \texttt{85b6f2d} & reply 線形性と期限の検査 \\
aios-claude & \texttt{9c7696c} & 期限検査がトップレベルを見ていなかった件 \\
aios-claude & \texttt{e7b6295} & select の返信経路 \\
\midrule
drone-hil & \texttt{b78ec34} & 拡張子を \texttt{.aipl} へ（3本）\\
drone-hil & \texttt{db7f5da} & JS-I を最新仕様へ追随 \\
drone-hil & \texttt{bee8408} & \verb|+| を数値専用に \\
drone-hil & \texttt{22c0569} & reply 線形性と期限の検査 \\
drone-hil & \texttt{c2929f5} & select の返信経路 \\
drone-hil & \texttt{4b6f2b5} & \texttt{wait} が \texttt{reply} を遅らせる \\
drone-hil & \texttt{601dd38} & 引数の型注釈 \\
\midrule
aipl\_line\_simulator & \texttt{7856575} & 拡張子を \texttt{.aipl} へ（2本）\\
\bottomrule
\end{longtable}
\end{center}

\noindent \texttt{abclcp} は push 済み。他の3リポジトリは作業中の feature
ブランチ上にあるためコミットのみで push は保留している。

\section{まとめ}

\begin{enumerate}
\item 拡張子を \texttt{.aipl} に統一した（4リポジトリ 478 本、参照 509 ファイル）。
  型検査の結果が改名前と 93/93 で一致することを確認している。

\item Py-I と JS-I を最新の言語仕様へ追随させた。ガイドのサンプルは
  \textbf{0/8 $\to$ 8/8}、しかも3実装で\textbf{出力が一致}する。

\item \verb|+| $\to$ \verb|++| の移行は\textbf{構文木で連鎖を判定}して行った
  （238本5448箇所）。正規表現では優先順位のせいで意味が変わる。

\item 揃える過程で\textbf{6つの不具合}が出た。うち3つは正であるはずの
  OCaml 版のものだった ---- トップレベル送信の順序破壊、
  アクターを引数で渡せない件（原因3つ）、そして \texttt{expr\_of\_value} が
  真偽値まで文字列に潰していたこと。

\item \texttt{g4\_select} の差を「スケジューリングの違い」で片付けず
  \emph{どちらが正しいか}を測ったことが、連鎖的な発見の起点になった。
  片付けていれば3つとも見つかっていない。

\item 見つけ方は3つの型に集約できる ----
  \textbf{3実装を並べて件数を比べる}、
  \textbf{機械的な置換をせず構文木で判定する}、
  \textbf{検査の指摘は件数でなく実物を開いて確かめる}。
  3つ目は 62 件の指摘のうち 59 件が誤検出であることを暴いた。
\end{enumerate}

\section*{参考}

\begin{itemize}
\item \texttt{docs/aipl\_guide\_ja.pdf} --- AIPL ユーザーズガイド（26ページ）。
  サンプル 8 本をソース全文と実行結果つきで解説。
\item \texttt{docs/aipl\_soundness2\_ja.pdf} --- 型健全性と型安全性の第2版。
  期限を必須にすると進行定理が二択になる。
\item \texttt{docs/aipl\_reply\_inference\_ja.pdf} ---
  \texttt{reply} からの戻り値型推論と注釈、線形性。
\item \texttt{docs/samples/guide/g1..g9} --- 本レポートで参照したサンプル。
\end{itemize}

\end{document}
